\begin{abstract}
Class-imbalance correction techniques -- resampling, cost-sensitive reweighting, and,
for tabular foundation models, test-time posterior adjustment -- are treated as
near-universally safe interventions: applying one is assumed to help, or at worst do
no harm. This assumption has been evaluated almost exclusively under the premise that
training labels are correct. We audit seven widely used correction methods, spanning
TabPFN (threshold correction, DistPFN, DistPFN-T, class-balanced downsampling) and
XGBoost (cost-sensitive weighting, random oversampling, SMOTE), under two directed
label-error mechanisms -- majority-to-minority contamination and minority-to-majority
loss -- across five UCI benchmark datasets, holding total training-set size fixed
across imbalance levels so that imbalance and label noise can be varied as independent
factors. Across 6{,}417 stress-tested evaluations, every method exhibits a nonzero
Correction Harm Rate: the fraction of conditions under which a correction
underperforms its uncorrected baseline ranges from 0.6\% to 31.7\% depending on method
and corruption direction, with break-even contamination rates as low as 5\% for some
method-dataset pairs. Friedman tests confirm these differences are statistically
significant ($p < 10^{-11}$) in every method family and corruption direction tested.
No single correction is robust across both corruption directions, and robustness
varies substantially by dataset. These results indicate that imbalance correction
cannot be assumed safe once label quality is uncertain, and that robustness to
directed label error should be evaluated alongside accuracy before deploying any
given correction in practice.
\end{abstract}
